\documentclass[a4paper,12pt]{article}

% --- PAQUETES ---
\usepackage[utf8]{inputenc}
\usepackage[spanish]{babel}
\usepackage{geometry}
\usepackage{setspace}
\usepackage{lmodern}
\usepackage{graphicx}
\usepackage{array}
\usepackage{booktabs}
\usepackage{titlesec}
\usepackage{fancyhdr}
\usepackage{xcolor}
\usepackage{colortbl}
\usepackage{enumitem}

% --- CONFIGURACIONES ---

\geometry{left=25mm, right=25mm, top=25mm, bottom=25mm}
\setstretch{1.2}

\definecolor{primary_blue}{RGB}{0, 84, 166}
\definecolor{dark_gray}{RGB}{64, 64, 64}
\definecolor{light_gray}{RGB}{240, 240, 240}

\pagestyle{fancy}
\fancyhf{}
\fancyhead[L]{\textcolor{dark_gray}{\textbf{WASAFF CONSULTING}}}
\fancyhead[R]{\textcolor{dark_gray}{Propuesta Técnica y Comercial}}
\fancyfoot[C]{\thepage}
\renewcommand{\headrulewidth}{0.4pt}
\renewcommand{\footrulewidth}{0.4pt}

\titleformat{\section}{\large\bfseries\color{primary_blue}}{\thesection}{1em}{}
\titlespacing*{\section}{0pt}{3.5ex plus 1ex minus .2ex}{2.3ex plus .2ex}

% --- DOCUMENTO ---
\begin{document}

% === PORTADA ===
\begin{titlepage}
    \centering
    \vspace*{1cm}
    
    {\scshape\huge \textbf{WASAFF CONSULTING}}
    \par\vspace{0.5cm}
    
    {\Large \textbf{Propuesta Técnica y Comercial}}
    \par\vspace{1cm}
    
    \color{dark_gray}{\rule{\linewidth}{1pt}}
    \par\vspace{0.7cm}
    
    \begin{tabular}{ll}
        \textbf{Fecha de Emisión:} & 22 de agosto de 2024 \\
        \textbf{Propuesta Nº:} & CLFWG-2024-001 \\
        \textbf{Validez:} & 20 días \\
    \end{tabular}
    
    \par\vspace{1cm}
    
    \color{dark_gray}{\rule{\linewidth}{0.4pt}}
    \par\vspace{0.4cm}
    \begin{tabular}{ll}
        \textbf{Preparado para:} & Sr. Francisco Pérez \\
        & Synergy Fire Protection \\
        \textbf{Referencia:} & Evaluación Predictiva de Riesgos Operacionales \\
    \end{tabular}
    \color{dark_gray}{\rule{\linewidth}{0.4pt}}
    
    \vfill 
    
    \begin{flushleft}
        \textbf{Estimado Sr. Pérez,}
        \vspace{0.5cm}
        
        Agradezco la oportunidad de presentar esta propuesta para un \textbf{Análisis de Ingeniería Predictiva sobre los Riesgos de Maniobras de Tuberías}.
        
        En un sector donde la seguridad es intransable, nuestro objetivo es claro: aplicar un análisis físico-matemático riguroso para transformar la incertidumbre operativa en certeza cuantificable. Este estudio le proporcionará una base técnica sólida para proteger a su personal, asegurar sus activos y fortalecer sus protocolos de seguridad ante sus clientes y entidades reguladoras.
        
    \end{flushleft}
\end{titlepage}

% === CUERPO DE LA PROPUESTA ===

\section*{1. Entendimiento del Desafío}

Synergy Fire Protection requiere una evaluación técnica independiente y de alto rigor para validar la seguridad en sus procedimientos de instalación de tuberías. El desafío principal es cuantificar los riesgos asociados a la posible caída de tuberías de acero de diversas dimensiones y pesos, con el fin de desarrollar protocolos de mitigación basados en evidencia, no en suposiciones.

\section*{2. Alcance y Metodología}

Para abordar este desafío, mi enfoque combina la física clásica con el análisis computacional:

\begin{enumerate}[leftmargin=*, label=\textbf{Fase \arabic*:}]
    \item \textbf{Modelado Físico-Matemático:} Se desarrollará un modelo que describe la dinámica de la caída de las tuberías (1'' a 12'', Sch 10-40) desde alturas variables (1 a 16 m), calculando parámetros críticos como la energía de impacto y el potencial de daño.
    \item \textbf{Análisis Paramétrico y Simulación:} Se ejecutarán simulaciones para un amplio rango de escenarios, identificando las combinaciones de peso, altura y trayectoria que presentan el mayor riesgo.
    \item \textbf{Desarrollo de Recomendaciones de Ingeniería:} Con base en los resultados cuantitativos, se formularán recomendaciones específicas para la mitigación de riesgos (ej. protocolos de izaje, zonas de exclusión, EPP recomendados, etc.).
\end{enumerate}

\section*{3. Entregables Clave}

Al finalizar el proyecto, usted recibirá un \textbf{Informe Técnico de Mitigación de Riesgos}, una herramienta estratégica que incluye:

\begin{itemize}[leftmargin=*, label=\textbullet]
    \item \textbf{Matriz de Riesgo Cuantificada:} Una clasificación clara de los escenarios analizados, desde riesgo bajo hasta crítico.
    \item \textbf{Resultados Visuales:} Gráficos y tablas que ilustran de manera inequívoca la relación entre las variables y el nivel de riesgo.
    \item \textbf{Protocolo de Recomendaciones Técnicas:} Un conjunto de directrices prácticas y fundamentadas para implementar en sus operaciones.
    \item \textbf{Reunión de Presentación y Asesoría:} Una sesión de trabajo para presentar los hallazgos y resolver todas las consultas de su equipo técnico.
\end{itemize}

\section*{4. Inversión}
La inversión para este servicio de consultoría especializada, que abarca la totalidad de los objetivos y entregables descritos, se presenta a continuación.

\begin{table}[h!]
    \centering
    \renewcommand{\arraystretch}{1.5}
    \begin{tabular}{|p{10cm}|c|}
        \hline
        \rowcolor{dark_gray!80}
        \textbf{\color{white}Descripción del Servicio} & \textbf{\color{white}Valor (UF)} \\
        \hline
        \rowcolor{light_gray}
        \textbf{Análisis Predictivo de Riesgos Operacionales} \newline \small{Servicio completo incluyendo modelado, simulación, informe y asesoría.} & 60.00 \\
        \hline
        \textbf{Subtotal} & \textbf{60.00} \\
        \hline
        \textbf{IVA (19\%)} & \textbf{11.40} \\
        \hline
        \rowcolor{light_gray}
        \textbf{Valor Total de la Inversión} & \textbf{71.40} \\
        \hline
    \end{tabular}
    \caption*{\small{\textit{El monto final en CLP se calculará según la UF del día de emisión de la factura.}}}
\end{table}

\section*{5. Términos y Condiciones}

\begin{itemize}[leftmargin=*, label=\textbullet]
    \item \textbf{Términos de Pago:} Se requiere un 50\% de la inversión (35.70 UF IVA incl.) para iniciar el proyecto. El 50% restante se abonará contra entrega satisfactoria del informe y la reunión de presentación.
    \item \textbf{Cronograma:} El plazo estimado de entrega es de 5 días hábiles a partir de la recepción del pago inicial y de toda la información técnica necesaria por parte del cliente.
    \item \textbf{Confidencialidad:} Toda la información compartida por Synergy Fire Protection será tratada con la más estricta confidencialidad.
\end{itemize}

\vfill 

Agradezco su confianza y quedo a su entera disposición para discutir esta propuesta en mayor detalle.

Atentamente,
\vspace{1.5cm}

\noindent\rule{7cm}{0.4pt}\\
\textbf{Felipe Wasaff González}\\
Físico Computacional y Consultor Principal\\
Wasaff Consulting\\
felipe@wasaffconsulting.com | +56 9 4612 5682

\end{document}